\subsection{Comentarii moderne}

\subsubsection{Perspectiva Ortodoxă}
% Idee extrasă din comentarii biblice ortodoxe moderne
% Exemplu: Bartolomeu Anania, Pr. Ioan Sorin Usca
\begin{displayquote}
    „Desigur că pentru autorul Genezei, lucrul principal nu este fructul din care au mâncat oamenii, și nici urmarea acestei acțiuni.
    Lucrul acesta se vede din versetele 14-17 ale capitolului 3, unde omul
    și femeia sunt blestemați pentru neascultare. După  capitolul 3 al
    genezei, păcatul nu este altceva decat neascultare de porunca divină”
    (Pr. Prof. Athanase Negoiță,
    \textit{Teologia Biblică a Vechiului Testament},
    Editura Sophia,
    București, 2004, p. 151)
\end{displayquote}

\subsubsection{Perspectiva Eterodoxă (receptată critic)}
% Idee extrasă din comentarii biblice catolice sau protestante
% Exemplu: Word Biblical Commentary, Jerome Biblical Commentary
\begin{displayquote}
    „Dacă el prin ascultarea poruncii dumnezeeşti ar fi rămas în comunitate
    cu Dumnezeu, atunci ar fi putut mânca din pomul acesta, fiindcă eră creat
    spre o vieaţă de veci. ” (Sfântul Ioan Casian, \textit{Convorbiri duhovnicești} 2, 3)
\end{displayquote}

\subsection{Mini-eseu: Importanța teologică a termenului}
% Aici scrii jumătatea de pagină de comentarii proprii (aproximativ 250-300 cuvinte)
% Eseul trebuie să surprindă importanța teologică a termenului în cele șapte versete selectate.
Textul eseului tău începe aici...
Acesta va analiza modul în care versetele din Vechiul Testament, împreună cu cele din Noul Testament, construiesc un înțeles teologic coerent al termenului [TERMENUL TĂU], prin prisma sfinților părinți și a exegezei moderne.
